% !TEX root = ../diary-book.tex

\section{令和7年08月23日 空気シャワー粒子の到来頻度とセットアップ}
\label{sec:diary-20250823}

空気シャワー粒子は主に、ガンマ線、電子、ミューオン（とニュートリノ）である。ニュートリノは検出が難しいため、ここではガンマ線、電子、ミューオンに注目する。
ガンマ線の検出効率は場合によっては低いこととなってしまうが、検出器への要請で頑張ることにして、ひとまず検出効率 100\% が達成できるとする。

サイトは 100~m 四方に 1~m$^2$のサイズの検出器をおいたものとする。検出器は 10~m 間隔で 11 x 11 = 121 台設置してあることになる。
各サイトでエネルギー損失の合計（時間の同時計数は見ていない）が 0.01~GeV 以上もつものがヒットしたサイトとする。各サイトにおいて検出されたガンマ線、電子、ミューオンの個数、エネルギー分布を調べる。


ミューオンの場合は中心から大きく離れた成分も一つのヒットとして使えるが、シャワーの場合は芯しかつかえないので、効率が悪い可能性が高い。なるほど。